\chapter{Estado da Arte}
\label{chap:estado-arte}

\section{Introdução}
\label{sec:ea-intro}

Este capítulo apresenta o levantamento do estado da arte nas áreas fundamentais para o desenvolvimento deste projeto: gestão financeira pessoal (\ac{PFM}), Open Banking e a aplicação de \ac{ML} no setor financeiro. Adicionalmente, é realizada uma análise comparativa detalhada dos principais concorrentes, tanto a nível internacional como europeu e português.

\section{Gestão Financeira Pessoal}
\label{sec:ea-pfm}

O conceito de \ac{PFM} refere-se ao conjunto de ferramentas e práticas que permitem aos indivíduos monitorizar, gerir e planear as suas finanças pessoais~\cite{pfm2023}. Historicamente, a gestão financeira era realizada através de folhas de cálculo ou registos manuais. Com a digitalização, surgiram plataformas especializadas que automatizam grande parte deste processo.

As plataformas modernas de \ac{PFM} oferecem tipicamente as seguintes funcionalidades: agregação de contas bancárias, categorização automática de transações, criação de orçamentos, definição de metas de poupança e geração de relatórios financeiros. Algumas soluções mais avançadas incluem ainda investimentos integrados, simulação fiscal e aconselhamento financeiro com recurso a \ac{IA}~\cite{fintech2024}.

\subsection{Evolução Histórica}

A primeira geração de ferramentas \ac{PFM} incluiu soluções como o Microsoft Money (1991) e o Quicken (1983), que funcionavam como software desktop. A segunda geração, liderada pelo Mint (2006), introduziu a agregação bancária online e a categorização automática de transações. A geração atual é caracterizada pela integração de \ac{IA}, Open Banking regulamentado e experiências mobile-first~\cite{pfm-evolution2024}.

\subsection{O Mercado Português}

Em Portugal, o panorama de \ac{PFM} é relativamente limitado. O Moey!, pertencente ao grupo Crédito Agrícola, é o principal neobank português, oferecendo funcionalidades básicas de categorização e objetivos de poupança (``Potes''). Contudo, não suporta agregação de contas externas via Open Banking, nem oferece funcionalidades avançadas de \ac{IA} ou simulação fiscal. Os neobancos europeus presentes em Portugal (Revolut, N26, Bunq) oferecem analytics financeiros, mas sem qualquer adaptação ao contexto fiscal e bancário português.


\section{Open Banking e PSD2}
\label{sec:ea-openbanking}

\subsection{A Diretiva PSD2}

A \ac{PSD2} (Payment Services Directive 2) é uma diretiva da União Europeia que entrou em vigor em janeiro de 2018, regulamentando os serviços de pagamento no mercado europeu~\cite{psd2directive}. Um dos seus principais contributos foi a obrigatoriedade de os bancos disponibilizarem \ac{API}s abertas que permitem a terceiros autorizados aceder à informação das contas dos clientes (com o seu consentimento explícito).

A diretiva define dois tipos principais de prestadores de serviços de terceiros:
\begin{itemize}
    \item \textbf{\ac{AISP} (Account Information Service Provider):} Entidades autorizadas a aceder à informação das contas bancárias para fins de agregação e análise.
    \item \textbf{PISP (Payment Initiation Service Provider):} Entidades autorizadas a iniciar pagamentos em nome do cliente.
\end{itemize}

\subsection{Implementação em Portugal}

Em Portugal, a implementação da \ac{PSD2} é supervisionada pelo Banco de Portugal. A \ac{SIBS}, enquanto operadora da rede Multibanco, disponibiliza a SIBS API Market, que permite o acesso a serviços de Open Banking de bancos portugueses. Adicionalmente, agregadores europeus como a Salt Edge oferecem cobertura de bancos portugueses, incluindo a Caixa Geral de Depósitos, o Millennium BCP, o Santander Portugal e o BPI~\cite{saltedge2024}.

\subsection{Salt Edge como Agregador}

A Salt Edge é uma plataforma de agregação financeira europeia que oferece uma \ac{API} \ac{REST} para acesso a contas bancárias via \ac{PSD2}. A plataforma disponibiliza um ambiente sandbox gratuito para desenvolvimento e testes, com cobertura de mais de 5.000 instituições financeiras em 50 países, incluindo os principais bancos portugueses~\cite{saltedge2024}. Para o contexto deste projeto, a Salt Edge foi selecionada como o agregador principal devido à sua cobertura do mercado português, à disponibilidade de sandbox gratuita e à qualidade da documentação técnica.


\section{Machine Learning em Fintech}
\label{sec:ea-ml}

A aplicação de \ac{ML} no setor financeiro (fintech) tem crescido exponencialmente nos últimos anos, abrangendo áreas como a deteção de fraude, a categorização de transações, a avaliação de risco de crédito e o aconselhamento financeiro automatizado (robo-advisors)~\cite{ml-fintech2024}.

\subsection{Categorização Automática de Transações}

A categorização de transações bancárias é um problema clássico de classificação de texto em \ac{ML}. As transações bancárias contêm tipicamente uma descrição textual curta (e.g., ``COMPRA CONTINENTE COVILHA''), um montante e uma data. O objetivo é atribuir automaticamente cada transação a uma categoria pré-definida (e.g., ``Supermercado'', ``Transportes'', ``Saúde'').

As abordagens mais comuns incluem:
\begin{itemize}
    \item \textbf{Modelos baseados em regras:} Utilizam expressões regulares e dicionários de comerciantes para classificação determinística.
    \item \textbf{Modelos de \ac{ML} tradicionais:} Algoritmos como Naive Bayes, SVM (Support Vector Machine) e Random Forest aplicados a features extraídas do texto das transações~\cite{sklearn2024}.
    \item \textbf{Modelos de Deep Learning:} Redes neuronais recorrentes (LSTM) ou transformadores (BERT) para classificação de sequências textuais.
\end{itemize}

Para o contexto deste projeto, optou-se pela utilização da biblioteca scikit-learn (Python), combinando técnicas de TF-IDF para vetorização do texto com um classificador Random Forest, por oferecer um bom equilíbrio entre precisão e complexidade de implementação no contexto de um projeto académico.

\subsection{IA Generativa em Finanças}

A \ac{IA} generativa, nomeadamente modelos de linguagem como o GPT-4o da OpenAI, tem encontrado aplicações no setor financeiro para geração de insights personalizados, chatbots de aconselhamento financeiro e análise de sentimento de mercado~\cite{openai2024}. Contudo, a aplicação de \ac{LLM}s em domínios com requisitos de precisão numérica e legal elevada --- como o cálculo fiscal --- levanta questões críticas de fiabilidade, auditabilidade e propensão para erros factuais (\textit{hallucinations})~\cite{llm-reliability2024}. Por este motivo, o presente projeto opta por uma abordagem alternativa para o módulo fiscal, descrita na secção seguinte.


\subsection{IA em Otimização Fiscal}
\label{sec:ea-fiscal-ia}

A aplicação de \ac{IA} à otimização fiscal constitui uma área de investigação crescente no domínio da fintech. Ao contrário da categorização de transações --- onde uma classificação aproximada é funcionalmente aceitável --- o cálculo fiscal exige precisão absoluta e rastreabilidade completa, uma vez que os resultados têm implicações legais e financeiras diretas para o utilizador~\cite{tax-ai2023}.

As abordagens existentes na literatura dividem-se em três paradigmas principais:

\begin{itemize}
    \item \textbf{Sistemas baseados em regras (\textit{expert systems}):} Codificam diretamente as disposições do código fiscal como lógica determinística. Garantem precisão total e auditabilidade, mas requerem atualização manual a cada alteração legislativa.
    \item \textbf{Modelos de \ac{ML} para classificação fiscal:} Utilizam algoritmos supervisionados para classificar transações segundo a sua dedutibilidade (saúde, educação, habitação), com base em padrões de texto e montante. São especialmente eficazes quando treinados com dados derivados diretamente das regras do código fiscal.
    \item \textbf{Algoritmos de otimização:} Dado um conjunto de rendimentos e deduções identificadas, aplicam técnicas de otimização combinatória (\textit{greedy}, \textit{branch-and-bound}) para determinar o cenário de declaração mais vantajoso entre as alternativas possíveis (e.g.\ tributação individual vs.\ conjunta).
\end{itemize}

A literatura académica recente demonstra que a abordagem híbrida --- \ac{ML} para identificação de despesas dedutíveis, motor determinístico para cálculo fiscal e algoritmo de otimização para comparação de cenários --- oferece o melhor equilíbrio entre adaptabilidade e precisão em sistemas de aconselhamento fiscal automatizado~\cite{hybrid-tax-ai2024}. Esta é a abordagem adotada no presente projeto, tendo como fontes de dados fidedignas e oficiais o Código do IRS português (\ac{CIRS} 2025) e o Orçamento do Estado 2025 (Lei n.º~24-D/2024, publicada no Diário da República).

A decisão de não recorrer a \ac{LLM}s externos para o módulo de otimização fiscal é intencional e assenta em três razões fundamentais: (1)~\textbf{falta de auditabilidade} --- os modelos generativos não garantem que os cálculos respeitem as regras exatas do \ac{CIRS}; (2)~\textbf{imprecisão numérica} --- \ac{LLM}s são reconhecidamente propensos a erros em operações aritméticas complexas; (3)~\textbf{dependência externa} --- a fiabilidade de uma funcionalidade crítica não deve depender da disponibilidade de \ac{API}s de terceiros. A IA desenvolvida é totalmente própria, treinada com dados oficiais e executada localmente na infraestrutura do projeto.


\section{Análise Comparativa de Concorrentes}
\label{sec:ea-concorrentes}

\subsection{Concorrentes Internacionais}

Foram analisados os quatro principais concorrentes no mercado internacional de \ac{PFM}:

\subsubsection{Copilot Money}
O Copilot Money é uma aplicação exclusiva para iOS que se destaca pelo design premium (finalista do Apple Design Award 2024) e pela precisão da sua categorização com \ac{ML}, com apenas cerca de 20\% das transações a necessitarem de recategorização. No entanto, é exclusivo da plataforma Apple, não suporta funcionalidades para casais e tem um preço de \$95/ano~\cite{copilot2024}.

\subsubsection{Monarch Money}
O Monarch Money é uma plataforma multiplataforma (web, iOS, Android) com foco em funcionalidades familiares, incluindo contas separadas e conjuntas para casais. Integra um assistente de \ac{IA} baseado em GPT-4, dashboards visuais avançados e um recurso de ``Month in Review'' automático. Suporta mais de 13.000 instituições financeiras, com um preço de \$99,99/ano~\cite{monarch2024}.

\subsubsection{Treasury}
O Treasury diferencia-se pela máxima flexibilidade, oferecendo um sistema de ``blocos de construção'' que permite ao utilizador criar o seu próprio sistema de orçamentação. Esta abordagem modular é especialmente atrativa para utilizadores avançados, embora possa apresentar uma curva de aprendizagem mais acentuada para iniciantes.

\subsubsection{Origin Financial}
O Origin Financial é a solução mais completa (``all-in-one''), integrando orçamentação, investimentos, declaração de impostos e planeamento patrimonial. Destaca-se pela conta poupança com 4,37\% APY e por regras avançadas de categorização automática, com um preço de \$99/ano.


\subsection{Concorrentes Europeus e Portugueses}

\subsubsection{Revolut}
Com mais de 70 milhões de utilizadores, o Revolut é o maior neobank europeu. Oferece suporte para 34+ moedas, investimentos (ações e criptomoedas), analytics de gastos e ligação a contas externas. No entanto, não tem foco no mercado português, não integra com o sistema fiscal e a interface pode ser complexa para utilizadores menos técnicos.

\subsubsection{N26}
O N26 é um banco digital alemão com licença bancária completa, oferecendo proteção de depósitos até 100.000 euros. Destaca-se pela interface minimalista e onboarding rápido, mas tem funcionalidades de analytics e Open Banking limitadas.

\subsubsection{Moey!}
O Moey! é o principal neobank português, pertencente ao grupo Crédito Agrícola. É totalmente gratuito e integra nativamente com MB Way e a rede Multibanco. Contudo, apresenta funcionalidades muito básicas de analytics, não suporta Open Banking para contas externas e não tem funcionalidades de \ac{IA} ou simulação fiscal.


\subsection{Análise Comparativa}

A Tabela~\ref{tab:comparativa} sintetiza a análise comparativa das principais plataformas analisadas.

\begin{table}[H]
\centering
\caption{Análise comparativa dos concorrentes principais.}
\label{tab:comparativa}
\small
\begin{tabularx}{\textwidth}{l c c c c c c}
\toprule
\textbf{Característica} & \textbf{Copilot} & \textbf{Monarch} & \textbf{Origin} & \textbf{Revolut} & \textbf{Moey!} & \textbf{Gold Lock} \\
\midrule
Multi-plataforma         & --  & Sim & Sim & Sim & Sim & Sim \\
IA/ML                    & Sim & Sim & Sim & --  & --  & Sim \\
Open Banking PT          & --  & --  & --  & Sim & --  & Sim \\
Simulador IRS            & --  & --  & Parcial & -- & -- & Sim \\
IA Fiscal Própria        & --  & --  & --  & --  & --  & \textbf{Sim} \\
Integração Seg. Social   & --  & --  & --  & --  & --  & \textbf{Sim} \\
Casais/Família           & --  & Sim & Sim & --  & --  & Futuro \\
Categorias PT            & --  & --  & --  & --  & Básico & Sim \\
Design Liquid Glass      & --  & --  & --  & --  & --  & Sim \\
Preço                    & \$95 & \$100 & \$99 & Grátis & Grátis & TBD \\
\bottomrule
\end{tabularx}
\end{table}


\section{Lacunas Identificadas no Mercado}
\label{sec:ea-lacunas}

A análise realizada permitiu identificar as seguintes lacunas críticas no mercado, que nenhum concorrente atual colmata:

\begin{enumerate}
    \item \textbf{Integração com o sistema fiscal português:} Nenhuma plataforma oferece simulação de \ac{IRS} ou integração com a \ac{AT}/Portal das Finanças.
    \item \textbf{IA fiscal própria e auditável:} Nenhum concorrente oferece um sistema de \ac{IA} desenvolvido internamente para otimização de \ac{IRS}, que cruze transações bancárias com dados da Segurança Social e as regras do \ac{CIRS}, sem depender de \ac{LLM}s externos.
    \item \textbf{Categorização adaptada ao mercado português:} As categorias de gastos não contemplam especificidades portuguesas (portagens, CTT, supermercados nacionais).
    \item \textbf{Open Banking otimizado para Portugal:} Nenhuma plataforma é nativamente otimizada para a rede \ac{SIBS}/Multibanco.
    \item \textbf{Design Liquid Glass em fintech:} Nenhum concorrente adota o paradigma de design Liquid Glass.
\end{enumerate}

Estas lacunas representam a oportunidade que o presente projeto se propõe explorar.


\section{Conclusões}
\label{sec:ea-conclusoes}

O levantamento do estado da arte confirmou a viabilidade técnica e a relevância do projeto proposto. As tecnologias necessárias --- Open Banking via \ac{PSD2}, \ac{ML} para categorização e classificação fiscal de transações, e frameworks modernos de desenvolvimento web --- estão maduras e acessíveis. Simultaneamente, a análise comparativa revelou uma clara lacuna no mercado para uma plataforma de \ac{PFM} focada no mercado português, que combine funcionalidades avançadas de \ac{IA} com adaptação profunda ao contexto fiscal e bancário nacional. Em particular, a inexistência de qualquer solução que ofereça \ac{IA} de otimização fiscal própria --- baseada no \ac{CIRS} e integrada com dados da Segurança Social --- representa a oportunidade diferenciadora central deste projeto.
